% !TeX root = Master.tex
\section{Introduction}
\label{sec:introduction}
It is common for a researcher working on an applied study to struggle with the choice of a statistical method. For prediction, models can be tested by setting aside part of the initial sample as a test set and comparing predictions against observed outcomes. However, much current scientific research, particularly in fields such as medicine and economics, seeks to answer causal questions, often requiring the estimation of counterfactuals, which are unobserved by definition. Therefore, causal estimates cannot be evaluated against an observed ground truth.

Theoretical properties can guide this choice, but asymptotic approximations may be inaccurate in small samples. Also, relationships between variables may be complex, and non-ideal scenarios such as limited overlap or confounding are common. To circumvent this issue, researchers have tested these estimators on constructed datasets in controlled scenarios with a specified reference value for the estimand. By comparing estimates to this value, they can obtain an estimation error per dataset, and by aggregating and averaging over datasets, they can compute metrics that allow researchers to rank the estimators according to their performance.

This general strategy is known as a \emph{Monte Carlo benchmark study}. The degree to which conclusions drawn from such studies can be useful in a particular real-world setting depends on how well the simulation mechanism reproduces the features of the target data-generating process (DGP) that determine estimator performance. To improve realism, empirical data are increasingly used to construct Monte Carlo datasets. In the econometrics literature, studies using this approach are referred to as \emph{empirical Monte Carlo studies} (EMCS). \authorcite{advaniMostlyHarmlessSimulations2019} distinguish \emph{placebo designs}, in which the treated sample is selected from an empirical control group using an artificial propensity score, so causal effects are known to be null, from \emph{structured designs}, in which the empirical DGP is modeled and data are generated from that model to approximate the original process. In medicine and epidemiology, the related trend has been associated with \emph{statistical plasmodes}, which can be seen as a subset of EMCS in which covariates are resampled from empirical data and outcomes are generated using an outcome-generating model \cite{schreckStatisticalPlasmodeSimulations2024}.

In this paper, we propose a framework that formalizes an EMCS benchmark as a pair consisting of two distinct objects. The first is a dataset-level distribution from which full simulated datasets are generated, intended to approximate the target redrawing scheme under the unknown distribution that generated the observed data. This law determines the ideal sampling behavior of the estimators. The second is a benchmark value for the estimand of interest, against which estimates are compared. In Section~\ref{sec:framework}, we make explicit the assumptions underlying this construction and define the conditions under which a constructed benchmark can be considered useful for estimator comparison.

Traditionally, studies benchmarked estimators (often newly proposed ones) under user-specified DGPs. These DGPs have been criticized as simplistic or unrealistic~\supercite{huberPerformanceEstimatorsBased2013,gentzelCaseEvaluatingCausal,boulesteixIntroductionStatisticalSimulations2020}. In this study, we focus on structured EMCS, which aim to approximate the DGP behind a reference dataset, providing context-specific information about estimator performance that may help reduce model uncertainty. From now on, EMCS refers to structured EMCS. In response to these criticisms, flexible generative models have gained attention. For example, to avoid researcher-specified DGPs that could be suspected of favoring particular estimators, \authorcite{atheyUsingWassersteinGenerative2024} train Wasserstein generative adversarial networks (WGANs) to generate artificial data closely resembling a reference dataset. To support the resulting benchmark, they present diagnostics comparing the synthetic and original data. More recent approaches, including RealCause~\supercite{nealRealCauseRealisticCausal2021}, Credence~\supercite{parikhValidatingCausalInference}, and Frugal Flows~\supercite{manelaMarginalCausalFlows2024}, combine flexible generative models with control over causal characteristics of the simulated data. Together, these approaches illustrate a broader shift from researcher-specified simulation models toward data-adaptive generative mechanisms aiming for greater realism.

This approach also has limitations. \authorcite{advaniMostlyHarmlessSimulations2019} theoretically and empirically show that incorrectly modeling the original DGP can result in simulated datasets that bias some of the estimators, possibly pushing the researcher to choose a candidate that in reality underperforms when applied to data coming from the real DGP. \authorcite{gentzelCaseEvaluatingCausal} address the issue indirectly when discussing high ranking variability when changing the simulation model. More recently, \authorcite{amaranathImprovingGenerativeMethods} show that different generative methods and parameter settings can yield different assessments of causal estimator performance. They propose a simulation-based inference framework that weights these methods and settings according to their compatibility with the source data. \authorcite{shawCautionaryNotePlasmode2026} show that plasmode simulations can bias propensity score--based estimators when covariates and treatment are bootstrapped together. We will discuss this issue in more depth with regard to our estimand in Section~\ref{sec:framework}.

A particular concern is what we call \emph{model alignment bias}: the distortion in assessing estimator performance when the mechanism used to generate Monte Carlo datasets (the simulator) favors estimators that share one or more of its model specifications. In this paper, we focus on the model alignment bias that can arise from outcome modeling. Consider the following hypothetical situation:

\begin{enumerate}
    \item A researcher has an observed dataset and wishes to determine whether an estimator that is based on a random forest (RF) outcome model or one based on a gradient boosting (GB) outcome model is more appropriate for estimating a causal effect. The criterion is having the lowest mean squared error (MSE) under a particular redrawing scheme.

    \item If the researcher knew the original data-generating process (DGP), it would be possible to generate datasets, apply both estimators, and compare estimates against the true causal effect. Suppose that, in this hypothetical benchmark, RF would outperform GB.

    \item Since the original DGP is unknown, the researcher instead sets up an EMCS that fits a GB outcome model to the observed data and uses it to simulate Monte Carlo datasets and to specify a corresponding causal reference value. Both estimators are applied to the simulated datasets, and their MSEs are computed relative to the simulator's causal reference value.

    \item Due to the alignment between the GB simulator and the GB-based estimator, GB may achieve a lower MSE than RF under the simulated law, reversing their ranking under the hypothetical ideal benchmark. The researcher then concludes that GB is the better estimator, even though RF would have been preferred under the real DGP.
\end{enumerate}

The potential consequences of model alignment have also been investigated in the literature. \authorcite{schulerSynthValidationSelectingBest2017} find some evidence for its effect in linear regression and for an unconstrained gradient-boosting simulator, but not for their constrained gradient boosting approach that enforces a synthetic treatment effect, suggesting that the magnitude of its influence can depend on the model. \authorcite{karimPlasmodeBenchmarkCan2026} studies the limitations of plasmode benchmarking through three levels of transfer from the simulated to the source DGP: estimand transfer, per-estimator performance transfer, and ranking transfer. The author shows the existence of model alignment bias, under the notion of \emph{projection}, in the case of a parametric logistic outcome model both theoretically and with simulations.

In this work, we investigate the ability of fitted outcome simulators to reproduce causal estimator rankings under the true data-generating process. We consider 12 simulation settings, drawn from the ACIC 2017 designs \cite{hahnAtlanticCausalInference2019} and those of \authorcite{nieQuasioracleEstimationHeterogeneous2021}, with varying outcome and treatment-assignment mechanisms, error structures, and sample sizes to assess simulators and estimators across a range of conditions. While much of the existing literature considers population-level treatment effects and simulation designs that redraw covariates and treatment assignments, we focus on the sample average conditional treatment effect (SACTE), holding covariates and treatment assignments fixed across Monte Carlo replications. This allows us to isolate the effects of outcome-model approximation conditional on the observed covariates and treatments. We evaluate five outcome simulators, from parametric regression to flexible machine-learning methods, by comparing the estimator rankings they induce with those obtained under the true DGP. We refer to agreement between these rankings as ranking fidelity. We consider 15 causal estimators, deliberately including several that employ the same outcome-model families as the simulators to help assess the impact of model alignment. Since the conditional-mean learner underlying a simulator can also be used for direct effect estimation, we study whether using it to select among estimators adds value. We also compare ranking fidelity with estimand fidelity. We examine whether averaging the rankings induced by the five simulators improves ranking fidelity and whether marginal outcome-distribution distances help compare simulators.

The rest of the paper proceeds as follows. Section~\ref{sec:framework} introduces a general framework for EMCS in causal inference, discusses the choice of estimand, clarifies the distinction between the target DGP and the simulator-induced benchmark, and formalizes the notion of simulator fidelity, including how outcome-model approximation can reverse estimator rankings. Sections~\ref{sec:simulation-design}--\ref{sec:results} then present the simulation design choices, algorithmic workflow, and results. Finally, Section~\ref{sec:discussion} concludes with a discussion of the main implications and limitations.
